\documentclass[10pt,a4paper]{article}
\usepackage{xeCJK}
\usepackage{fontspec}
\usepackage{booktabs}
\usepackage{array}
\usepackage{xcolor}
\usepackage{geometry}
\usepackage[protrusion=true]{microtype}
\geometry{margin=1.7cm}
\linespread{1.05}

\setmainfont{Baskerville}
\setCJKmainfont{Songti SC}
\newfontfamily{\starfont}{Songti SC}

\definecolor{wrongred}{RGB}{198,40,40}
\definecolor{rulegray}{RGB}{190,190,190}

\newcommand{\checkmarkbox}{\fbox{\rule{0pt}{0.75ex}\rule{0.75ex}{0pt}}}
\newcommand{\num}[1]{\makebox[1.9em][r]{#1.}\hspace{0.3em}}
\newcommand{\blank}{\underline{\hspace{2.1cm}}}
% 判错的条：题号后加红星
\newcommand{\STAR}{\textcolor{wrongred}{\starfont ★}}
% ★ 固定宽度槽位：带星/不带星行正文左边缘一致
\newcommand{\STARSLOT}[1]{\makebox[1.35em][l]{#1}}
\newcommand{\zhs}[2]{\hangindent=2.2em\hangafter=1\num{#1}\STARSLOT{\STAR}#2}
\newcommand{\zhx}[2]{\hangindent=2.2em\hangafter=1\num{#1}\STARSLOT{}#2}
% 表格正文列：左对齐，避免 p 列两端对齐造成的 Underfull
\newcolumntype{L}[1]{>{\raggedright\arraybackslash}p{#1}}

\begin{document}
\pagestyle{plain}

\begin{center}
{\LARGE\bfseries 英语短语默写 \quad 二刷（中 $\to$ 英）}\\[0.25em]
{\large 小蓝 P21--30}\\[0.22em]
{\small 2029届高一英语 · \STAR 为 2026-10 原卷判错的 16 条}
\end{center}

\vspace{0.35em}
\noindent\begin{tabular}{@{}p{0.33\textwidth}@{}p{0.34\textwidth}@{}p{0.33\textwidth}@{}}
姓名：\underline{\hspace{2.2cm}} & 日期：\underline{\hspace{2.2cm}} & 得分：\underline{\hspace{1.5cm}}\,/\enspace 25 \\
\end{tabular}

\vspace{0.45em}
\noindent{\small 说明：根据中文短语与提示词默写出英文短语，题号沿用原卷。\STAR\ 为原卷判错的条（16 条），本次需重点复现；无星号的条用于检验是否遗忘。判定与标准答案由 AI 给出，二刷前请对照原卷核对。}

\vspace{0.45em}
\renewcommand{\arraystretch}{1.5}
\noindent\begin{tabular}{@{}L{5.1cm}@{\hspace{0.25cm}}L{2.7cm}@{\hspace{0.5cm}}L{5.1cm}@{\hspace{0.25cm}}L{2.7cm}@{}}
\toprule
中文短语（提示词） & 英文默写 & 中文短语（提示词） & 英文默写 \\
\midrule
\zhx{1}{除了做某事别无选择（alternative）} & \blank & \zhs{14}{要求道歉} & \blank \\
\zhs{2}{业余自行车手} & \blank & \zhx{15}{为某事向某人道歉（apologize）} & \blank \\
\zhx{3}{…及其他东西（among）} & \blank & \zhs{16}{政府呼吁人人节约用水（appeal）} & \blank \\
\zhs{4}{拿游戏取乐（amuse）} & \blank & \zhs{17}{你有兴趣去合资公司上班吗？（appeal）} & \blank \\
\zhx{5}{数据分析} & \blank & \zhs{18}{下结论（arrive）} & \blank \\
\zhs{6}{对拖延症感到气愤} & \blank & \zhs{19}{请假（apply）} & \blank \\
\zhs{7}{对…负责（answer）} & \blank & \zhx{20}{派某人任某职} & \blank \\
\zhs{8}{一个又一个胜利} & \blank & \zhs{21}{如果…我将不胜感激（appreciate）} & \blank \\
\zhx{9}{带某人参观（around）} & \blank & \zhx{22}{对…欣赏（appreciation）} & \blank \\
\zhx{10}{预料明年的销售量增加} & \blank & \zhs{23}{不同意我辍学（approve）} & \blank \\
\zhs{11}{对…渴望（anxious）} & \blank & \zhx{24}{拱腰} & \blank \\
\zhs{12}{掌声四起} & \blank & \zhs{25}{安排医生给他看病（arrange）} & \blank \\
\zhs{13}{在我看来（appear）} & \blank & \\
\bottomrule
\end{tabular}

\vspace{0.5em}
\noindent{\color{rulegray}\rule{\textwidth}{0.5pt}}
\vspace{0.3em}
\noindent\textbf{复核记录}\hfill\small 完成默写后填写

\vspace{0.3em}
\noindent\begin{tabular}{@{}p{0.31\textwidth}@{}p{0.34\textwidth}@{}p{0.35\textwidth}@{}}
默写得分：\underline{\hspace{1.2cm}}\,/\,25
& 仍错误的题号：\underline{\hspace{3.2cm}}
& 下次复习日期：\underline{\hspace{2.0cm}} \\
\end{tabular}

\vspace{0.3em}
\noindent 本轮易错点：\checkmarkbox\ 介词搭配（to / for / of）\quad\checkmarkbox\ 固定搭配（appeal to / arrange for）\quad\checkmarkbox\ 漏词\quad\checkmarkbox\ 用词词形

\vspace{0.3em}
\noindent 复习状态：\checkmarkbox\ 已掌握\quad\checkmarkbox\ 需要巩固\quad\checkmarkbox\ 需重点复习

\end{document}
